% =============================================================================
%  CHAPTER 12 -- THE PHASE MAP: HOW THIS PROJECT WAS EXECUTED
% =============================================================================
\section{The phase map: how this project was executed}
\label{ch:phasemap}

Part~II is a record of what was actually done. Before the individual phases, this chapter
explains the execution model, because the model is part of the result: a project built in
thirteen sequential phases with a hard exit gate at each one produces different (better,
more honest) output than a project written in one pass.

\subsection{The rules the project ran under}
\label{sec:phasemap:rules}

From \code{PLAN.md}, in force for every phase:

\begin{enumerate}[label=\textbf{R\arabic*.}]
  \item \textbf{Golden rule.} No number appears in the report unless it was produced by
    checked-in code in this repository. Nothing is reported without a significance test,
    out-of-sample validation, and cost adjustment. A negative result, clearly explained,
    counts as success.
  \item \textbf{One phase per session.} Each session completes exactly one Part, commits
    it, and writes \code{HANDOFF.md} briefing the next session completely --- enough that
    a fresh agent with no memory can continue. This is why \code{HANDOFF.md} contains
    recovery commands rather than prose.
  \item \textbf{Exit gates.} Each Part lists preliminary checks that must pass before the
    Part is declared done (compiles, unit tests pass, every figure the report references
    exists and is produced by code, no raw data committed, \code{check.sh} green).
  \item \textbf{C++ is the simulation core; Python is research and plotting.} No strategy
    simulation logic lives only in Python.
  \item \textbf{Walk-forward, lookahead-free, gross and net everywhere.}
  \item \textbf{Significance or it didn't happen.} No claim without its test; no best-case
    number without its multiple-testing-corrected companion.
  \item \textbf{The report is built, not typed.} \code{report/build\_report.py}
    regenerates every figure and table from \code{results/} before rendering.
\end{enumerate}

\subsection{The environment, and how it shaped the design}
\label{sec:phasemap:env}

Every design decision in this repository is downstream of a constraint discovered in
Phase~\ref{ch:phase0} and recorded in \code{environment/env\_map.md}:

\begin{center}
\renewcommand{\arraystretch}{1.2}
\begin{tabular}{p{0.30\linewidth}p{0.16\linewidth}p{0.46\linewidth}}
\toprule
Resource & Status & Consequence for the design \\
\midrule
\code{g++} 12.2, \code{make} & available & the C++ engine builds with the system toolchain \\
\code{cmake} & installable from PyPI & \code{cmake -S src -B build}; \code{PATH} must include \code{\$HOME/.local/bin} \\
Python 3.11 + PyPI & reachable & numpy/pandas/scipy/statsmodels/matplotlib/sklearn all usable \\
\code{pip} & PEP-668 externally managed & every install needs \code{--user --break-system-packages} \\
github.com / codeload & reachable & the only in-sandbox route to real market data (repo tarballs) \\
Yahoo / stooq / Nasdaq APIs & \textbf{blocked} & the primary 16.7-year total-return panel had to be downloaded externally and committed as raw CSVs \\
apt / conda / Rust & blocked & no \code{texlive}, no Eigen, no system packages at all \\
\textbf{LaTeX engine} & \textbf{not installable} & \code{report.tex} is canonical LaTeX but cannot be compiled here; \code{build\_report.py} always renders a self-contained HTML preview with locally rasterised math \\
CPU / RAM & 2 cores, 3.8\,GB & datasets kept small (195 names, 4{,}190 bars); parallelism limited to \code{-j2} \\
\bottomrule
\end{tabular}
\end{center}

Two consequences deserve emphasis because they are unusual.

\paragraph{No Eigen, no LAPACK.} Linear algebra is hand-rolled in
\code{src/model/dense\_la.h}: an SPD/ridge Cholesky solver and a cyclic-Jacobi symmetric
eigen-solver (Section~\ref{sec:pca:jacobi}). This was a forced choice that turned out to
be a good one --- the code is a few hundred lines, cache-friendly for $n\le5$, and
unit-tested against reconstruction identities rather than against another library's
output.

\paragraph{No TeX engine.} The report cannot be compiled to PDF inside the sandbox. The
solution is a two-track build: \code{report/report.tex} (with \code{report/chapters/*.tex})
is written to compile cleanly anywhere a normal TeX distribution exists, and
\code{report/build\_report.py} additionally renders \code{report/report.html} --- a
self-contained book preview in which every display equation is rasterised locally by
matplotlib's \code{mathtext} and every inline formula is transliterated to HTML, so no
CDN, no network, and no MathJax is needed. That is why ``running
\code{python3 report/build\_report.py}'' is the single command that produces the readable
artifact.

\subsection{The thirteen phases at a glance}
\label{sec:phasemap:table}

\begin{center}
\renewcommand{\arraystretch}{1.25}
\small
\begin{tabular}{p{0.05\linewidth}p{0.26\linewidth}p{0.36\linewidth}p{0.22\linewidth}}
\toprule
Phase & Goal & Main deliverable & Chapter \\
\midrule
0 & Plan, environment recon, scaffolding & \code{PLAN.md}, \code{environment/}, directory skeleton, report skeleton & \ref{ch:phase0} \\
1 & Real data, documented universe & 195-name 2010--2026 total-return panel, cleaner, QA gate, data summary & \ref{ch:phase1} \\
2 & Statistical foundation toolkit & \code{scripts/statkit/}: ADF, KPSS, EG, Johansen, DW, LB, BP, White, HAC, OU, DM, bootstrap, Bonferroni, RC, SPA + 35 unit tests & \ref{ch:phase2} \\
3 & Dependency-free C++ engine skeleton & SoA buffer, incremental rolling regression, basket selector, signal, book, latency histograms, hand-rolled LA, 13 tests & \ref{ch:phase3} \\
4 & Correctness core & walk-forward loop, next-bar execution, cost model, no-lookahead test, brute-force cross-check, first real backtest & \ref{ch:phase4} \\
5 & Baseline model comparison & OLS / Ridge / Lasso / ElasticNet / PCR through one identical causal pipeline & \ref{ch:phase5} \\
6 & Statistical diagnostics & the full battery on the traded residuals: stationarity, cointegration, autocorrelation, heteroskedasticity, HAC before/after, half-life vs.\ holding period & \ref{ch:phase6} \\
7 & Portfolio book and risk controls & four concurrent strategies, inverse-vol with caps, VaR/CVaR, regime and stress decomposition, sensitivity heatmaps & \ref{ch:phase7} \\
8 & Kalman time-varying hedge & state-space model, Joseph-stabilised filter, synthetic validation, comparison vs.\ rolling OLS/Ridge & \ref{ch:phase8} \\
9 & Nonlinear extension + RESET & causal GBM/MLP forecasting, learning curves, RESET, gated-trading verdict & \ref{ch:phase9} \\
10 & Multiple testing + DM & hypothesis-count audit ($M{=}44$), Bonferroni, White RC, Hansen SPA, DM forest, power analysis & \ref{ch:phase10} \\
11 & Architecture and latency & design rationale, cache microbenchmarks, per-stage latency percentiles, traceability of every systems number & \ref{ch:phase11} \\
12 & Report assembly and non-goals & complete report, abstract, discussion, explicit non-goals, conclusion & \ref{ch:phase12} \\
13 & Reproducibility and delivery & end-to-end regeneration, this book, reproducibility manual, future work & \ref{ch:phase12} \\
\bottomrule
\end{tabular}
\end{center}

\subsection{Repository layout}
\label{sec:phasemap:layout}

\begin{verbatim}
environment/     bootstrap (setup.sh), exit gate (check.sh), env map, data candidates
scripts/
  data/          download.py, clean.py, qa.py, export_engine_input.py, sector map, tickers
  statkit/       the statistical library: stationarity, cointegration, autocorr,
                 heterosk, hac, ou, forecast (DM), bootstrap, multtest, misspec + tests/
  model/         baseline_models.py (walk-forward hedges + simulate), kalman.py,
                 nonlinear_models.py
  analysis/      capture_* (run the experiment, write results/tables/*.csv)
  plotting/      render_*  (read results/, write results/figures/*.png + report/tables/*.tex)
  test_kalman.py, test_nonlinear.py, run_statkit_tests.py
src/             the C++ engine (header-only)
  data/          csv_loader.h, market_data_buffer.h   (SoA ring buffer)
  model/         dense_la.h (Cholesky + Jacobi), rolling_regression.h, basket_selector.h
  signal/        signal_generator.h                   (z-threshold policy)
  portfolio/     portfolio_book.h, risk.h             (positions; inverse-vol + limits)
  engine/        backtest_engine.h, spread_backtest.h, latency_hist.h
  main.cpp       --version / --test / --demo
test/            minimal_test.hpp (hand-rolled harness) + engine_tests.cpp
bench/           microbench.cpp  (vector vs deque, SoA vs AoS)
report/          report.tex (root), chapters/*.tex (this book), references.bib,
                 build_report.py (the single entry point), tables/ + figures/ (generated)
data/            raw/yfinance/*.csv (committed), processed/ (git-ignored, regenerated)
results/         tables/*.csv, figures/*.png, metrics/ (all git-ignored, all regenerated)
PLAN.md          the 13-phase master plan and the rules above
HANDOFF.md       briefing for the next session, rewritten every phase
\end{verbatim}

Two conventions worth stating explicitly because they make the repository auditable:
\begin{itemize}
  \item \textbf{\code{capture\_*} versus \code{render\_*}.} Scripts named
    \code{capture\_\{backtest,portfolio,engine\_bench\}.py} or \code{compare\_*.py}
    \emph{run the experiment} and write machine-readable results to
    \code{results/tables/*.csv}. Scripts named \code{render\_*.py} \emph{only read}
    those CSVs and produce figures plus the \code{report/tables/*.tex} fragments. The
    separation means a rendering bug can never change a number, and a number can always
    be traced to one CSV and one capturing script.
  \item \textbf{Generated artifacts are git-ignored.} \code{results/},
    \code{data/processed/}, \code{build/}, \code{report/tables/},
    \code{report/figures/}, \code{report/report.html} are all reproducible and all
    excluded, so the repository holds sources only. The exception is
    \code{data/raw/yfinance/*.csv}, which is committed precisely because it
    \emph{cannot} be regenerated in this sandbox (Yahoo is blocked at runtime); it is the
    one piece of raw data that had to be frozen.
\end{itemize}

\subsection{How to reproduce everything, in one command sequence}
\label{sec:phasemap:reproduce}

\begin{verbatim}
export PATH="$HOME/.local/bin:$PATH"
bash environment/setup.sh                       # toolchain (~2 min, idempotent)
python3 scripts/data/clean.py                   # raw -> data/processed/universe.csv
python3 scripts/data/qa.py                      # 10/10 data-quality gate
python3 scripts/data/export_engine_input.py --target ADBE --basket CRM ADSK INTU
python3 scripts/data/export_engine_input.py --target AVGO --basket ADI INTC AAPL
python3 scripts/data/export_engine_input.py --target GS  --basket BAC
python3 scripts/data/export_engine_input.py --target CAT --basket HON
cmake -S src -B build && cmake --build build -j2
bash environment/check.sh                       # 52/52 PASS: every gate
python3 report/build_report.py                  # figures + tables + report.html
\end{verbatim}

Wall-clock on the sandbox (2 cores): setup $\approx2$ min, clean+QA $\approx40$ s, C++
build $\approx40$ s, \code{check.sh} (which runs every capture and render script plus all
test suites) $\approx7$ min, report build $\approx2$ min. Total $\approx12$ min from a
fresh clone to a finished book. Chapter~\ref{ch:reproducibility} documents the exact
expected outputs of each step so a run can be verified, not just executed.
